\clearpage\thispagestyle{empty}
\pdfbookmark[0]{Quellen zu Fadenintervalle (November 1823)}{wire1823-sources}
\begin{center}\textbf{Quellen zu Fadenintervalle (November 1823)}\end{center}
\noindent Diese KI-Übertragung ins heutige Deutsch folgt Gauß’ \textit{Fadenintervalle} in \href{https://gdz.sub.uni-goettingen.de/id/PPN236007467}{den \textit{Werken}, Band~VI} (Göttingen, 1874), S.~445--448. Der Aufsatz erschien zuerst als \textit{Neue Methode, die gegenseitigen Abstände der Fäden in Meridian-Fernröhren zu bestimmen} in den von H.~C. Schumacher herausgegebenen \textit{Astronomischen Nachrichten}, Band~II, Nr.~43 (November 1823), Sp.~\href{https://api.digitale-sammlungen.de/iiif/image/v2/bsb10049427_00218/full/1800,/0/default.jpg}{371--372}, \href{https://api.digitale-sammlungen.de/iiif/image/v2/bsb10049427_00219/full/1800,/0/default.jpg}{373--374} und \href{https://api.digitale-sammlungen.de/iiif/image/v2/bsb10049427_00220/full/1800,/0/default.jpg}{375--376}. Die Links führen zum Exemplar der Bayerischen Staatsbibliothek. Der Band erschien 1824 in Altona.
\par\smallskip\noindent Der Erstdruck nennt Gauß als \textit{Hofrath und Ritter} und trägt seine Unterschrift. Die Werkausgabe hat die Überschrift \textit{Fadenintervalle} und schreibt den Namen des Optikers \textit{Frauenhofer}, wo im Erstdruck \textit{Fraunhofer} steht. Beim zweiten wahren Abstand für Faden I am 23. Oktober steht in Sp.~374 des Erstdrucks 10 24,06. Die Übertragung behält die Lesung 10 24.86 der beschädigten Ziffer in der Werkausgabe bei. Die Einheitenzeichen stehen an den Stellen, an denen sie auch in der Werkausgabe stehen. \textit{Theoria Combinationis Observationum} ist die von Gauß verwendete Kurzform des Titels. KI-Übertragung und KI-Satz.
\clearpage
